\section{Introduction}

This semester's Assignment~3 requires using Vibe Coding---building software through natural language interaction with AI coding tools---to create a personal homepage.
The homepage must contain three required areas: self-introduction (name/major/year~$+$~a distinctive self-description~$+$~a representative image), data visualization (selecting a set of real personal data and presenting it in a suitable form), and an interactive easter egg (at least one triggerable interactive design).
Grading evaluates completeness, personalization, visual design and creativity, prompt engineering, and depth of reflection.

The core principle of Vibe Coding is ``humans steer the direction, AI executes the implementation.'' In practice, however, a non-technical user quickly faces a fundamental question: \textbf{how should I phrase my prompts?}
Vague descriptions like ``make me a nice page'' produce generic, impersonal results---the AI defaults to some template variant. Writing precise design descriptions, on the other hand, requires expertise in color theory, typography, and animation design.
How can a student with neither programming skills nor design background effectively steer AI in Vibe Coding? This is the question this experiment seeks to explore.

To find a Vibe Coding strategy suitable for non-expert users, this paper conducts a controlled self-experiment: first, using Codex to attempt \textit{Incremental Build}---having the AI build a page skeleton, then adding functional modules one at a time; then, switching to Claude Code to attempt \textit{Design-First}---investing substantial dialogue rounds in co-designing a complete specification with the AI before writing any code, compiling it into a design document, and implementing against that document.

The remainder of this paper is organized as follows: Section~2 briefly reviews the background of AI coding tools and prompt engineering. Section~3 details the two Vibe Coding strategies involved in this experiment. Section~4 presents the execution process and results of both strategies, including design philosophy, core interactions, the easter egg system, and the experience-driven iteration after the v1.0 implementation. Section~5 provides a summary and reflection.
